% QR-CODE: Lernpfad Induktion (Kl. 9, DS 09b)

\documentclass[a4paper,11pt]{article}

\usepackage[utf8]{inputenc}
\usepackage[T1]{fontenc}
\usepackage{helvet}
\renewcommand{\familydefault}{\sfdefault}

\usepackage[margin=20mm]{geometry}
\usepackage{tikz}
\usepackage{xcolor}
\usepackage{qrcode}

\definecolor{physik}{HTML}{2c5aa0}
\colorlet{fachfarbe}{physik}

\newcommand{\lptitel}{Elektromagnetische Induktion}
\newcommand{\lpurl}{https://devbox42.github.io/sciencesim/physik/kl09/motor-induktion/LP-09b-induktion.html}
\newcommand{\lpurlkurz}{devbox42.github.io/.../LP-09b-induktion.html}

\pagestyle{empty}

\begin{document}

\begin{center}

\vspace*{10mm}
{\fontsize{28}{34}\selectfont\bfseries\color{fachfarbe}Einführung: \lptitel}

\vspace{15mm}

\fcolorbox{fachfarbe}{white}{\qrcode[height=100mm]{\lpurl}}

\vspace{12mm}

{\large\color{gray}\lpurlkurz}

\vspace{8mm}

{\Large\color{gray}Scanne den QR-Code mit deinem Gerät}

\vspace{10mm}

\fcolorbox{fachfarbe}{fachfarbe!10}{%
    \parbox{0.7\textwidth}{%
        \centering
        \vspace{3mm}
        {\large Einführung zur elektromagnetischen Induktion}\\[2mm]
        6 Schritte -- ca. 12 Minuten
        \vspace{3mm}
    }%
}

\end{center}

\end{document}
